\section{Introducción: del juramento médico a la ética de la IA}

Esta clase es la novena sesión del curso de primavera de 2025 de MIT 6.S191 (Introduction to Deep Learning), impartida como invitado por Doug Blank, director de investigación de Comet ML. Doug Blank se dedica a la investigación de redes neuronales desde la década de 1990, completó su doctorado en ciencia cognitiva y ciencias de la computación en 1997, fue profesor durante 20 años en Bryn Mawr College y después se incorporó a Comet ML, una empresa emergente de gestión de experimentos de IA.

La pregunta central de esta clase es sumamente especulativa: \textbf{¿podemos hacer que los sistemas de IA hagan algo parecido al ``juramento hipocrático'' (Hippocratic Oath) de los médicos, es decir, comprometerse a ``no causar daño'' (Do No Harm)?}

\begin{importantbox}{Juramento hipocrático (Hippocratic Oath)}
El juramento hipocrático es un código de ética médica propuesto hace unos 2,500 años por Hipócrates, el padre de la medicina griega antigua. Antes de ejercer formalmente, los estudiantes de medicina juran cumplir con estándares éticos profesionales, cuyo principio central puede resumirse en una frase: \textbf{Do No Harm} (no causar daño). Esto expresa dos principios clave:
\begin{itemize}
    \item \textbf{Beneficence} (principio de beneficencia): buscar activamente el bienestar del paciente
    \item \textbf{Non-maleficence} (principio de no maleficencia): evitar causar daño al paciente
\end{itemize}
\end{importantbox}

El expositor señala que los problemas que enfrentan los sistemas de aprendizaje profundo van mucho más allá del plano técnico —la calidad de los datos de entrenamiento, el sobreajuste, la capacidad de generalización—, e incluyen también problemas sociales más amplios como el sesgo de los datos, el consumo de energía, la exactitud de las salidas y la seguridad. A medida que crecen la escala y la complejidad de los sistemas, estos problemas también se intensifican.

\begin{knowledgebox}{De 1997 a 2025: el salto de escala del aprendizaje profundo}
El expositor ilustra con su propia experiencia el cambio de escala del aprendizaje profundo: la red neuronal usada en su tesis doctoral de 1997 tenía apenas unas 30 neuronas y 3,000 parámetros, mientras que los modelos de lenguaje de gran escala actuales tienen cientos de miles de millones de parámetros. Los principios matemáticos son básicamente los mismos —multiplicación de matrices, funciones de activación, retropropagación—, pero la mejora en la eficiencia del hardware ha hecho que la escala de los sistemas crezca de forma exponencial, y con ello han surgido nuevos desafíos de ingeniería y de ética.
\end{knowledgebox}

\subsection{Resumen de la sección}

Esta clase se desarrolla en torno a una analogía central: si los profesionales de la medicina deben jurar ``no causar daño'', ¿deberían los profesionales de la IA hacer un compromiso similar? El expositor parte del marco de riesgo-beneficio propio de la inversión y, a partir de ahí, explora gradualmente la evaluación de riesgos de los proyectos de IA, la regulación normativa, las limitaciones técnicas y la responsabilidad de quienes desarrollan estos sistemas.
